% =====================================================================
%  FORMULIR PENGAJUAN JUDUL PENELITIAN / SKRIPSI / TUGAS AKHIR
%  Program Studi Rekayasa Sistem Komputer, Fakultas Teknik, UM Buton
%  Kompilasi PDF : xelatex (lihat menu.sh)   |   Word : pandoc (lihat menu.sh)
%  Susunan (5 hlm.): 1 identitas + tiga judul; 2-4 satu halaman per usulan; 5 hasil verifikasi.
%  Isi kotak: Arial 12 tanpa tebal. Petunjuk satu baris, miring.
% =====================================================================
\documentclass[12pt,a4paper]{article}
\usepackage{skripsi}
\usepackage[a4paper,top=1.5cm,bottom=1.5cm,left=2.5cm,right=2.5cm]{geometry}

% ---------- Data dosen & peran (common/dosen.tex, common/peran.tex) ----------
% =====================================================================
%  DAFTAR DOSEN  --  Program Studi Rekayasa Sistem Komputer, UM Buton
%  Satu dosen = dua makro:  \dsnKunci (nama dan gelar)  dan  \nidnKunci (NIDN).
%  Tambah dosen baru dengan menyalin dua baris di bawah. Peran (pembimbing, penguji, dst.)
%  diatur di peran.tex atau lewat  ./menu.sh  ->  "Atur dosen / peran".
% =====================================================================
\newcommand{\dsnAwaluddin}{Muhammad Awaluddin, S.Kom., M.Kom}
\newcommand{\nidnAwaluddin}{0909038201}
\newcommand{\dsnDarmawan}{Darmawan, S.Kom., M.Kom}
\newcommand{\nidnDarmawan}{0911048204}
\newcommand{\dsnRando}{Rando, S.Kom., M.Eng.}
\newcommand{\nidnRando}{0911069501}
\newcommand{\dsnSaidi}{La Ode Muhammad Saidi, S.Si., M.Kom}
\newcommand{\nidnSaidi}{0913038905}
\newcommand{\dsnDidin}{Didin Adri, S.Si., M.Sc}
\newcommand{\nidnDidin}{0930088802}
\newcommand{\dsnHibrian}{Muhammad Hibrian Wiwi, S.T., M.Kom}
\newcommand{\nidnHibrian}{0930129003}
\newcommand{\dsnRasyid}{Muhammad Abdur Rasyid Asruddin, S.T., M.Kom}
\newcommand{\nidnRasyid}{6641774675130242}
% Dekan Fakultas Teknik
\newcommand{\dsnAgusman}{Agusman, S.T., M.M}
\newcommand{\nidnAgusman}{0917088206}

% =====================================================================
%  PERAN DOSEN  --  dipakai otomatis di sampul, lembar persetujuan/pengesahan,
%  kata pengantar, dan abstrak semua dokumen.
%  Isi dengan kunci dosen dari dosen.tex, mis.  \dsnAwaluddin  dan  \nidnAwaluddin.
%  Atau jalankan  ./menu.sh  lalu pilih "Atur dosen / peran"  (berkas ini ditulis ulang otomatis).
% =====================================================================
\newcommand{\PembimbingI}{…………………}\newcommand{\PembimbingInidn}{……………………}
\newcommand{\PembimbingII}{…………………}\newcommand{\PembimbingIInidn}{……………………}
\newcommand{\PengujiI}{…………………}\newcommand{\PengujiInidn}{……………………}
\newcommand{\PengujiII}{…………………}\newcommand{\PengujiIInidn}{……………………}
\newcommand{\KetuaProdi}{…………………}\newcommand{\KetuaProdinidn}{……………………}
\newcommand{\Dekan}{\dsnAgusman}\newcommand{\Dekannidn}{\nidnAgusman}


% ---------- Makro yang ikut dibaca pandoc ----------
\newcommand{\logo}[1]{\includegraphics[width=#1]{../gambar/logo-umbuton.png}}
\newcommand{\titik}{……………………………………………}
\newcommand{\cek}{[\hspace{1.2em}]}
\newcommand{\tinggi}[1]{\rule{0pt}{#1cm}}

% Satu bagian isian: judul bagian (tebal), petunjuk satu baris (miring), ruang isi #3 cm
\newcommand{\bagian}[3]{{\raggedright\textbf{#1}\ {\footnotesize\itshape #2}\par}\tinggi{#3} \\ \hline}

% Satu usulan = satu halaman
\newcommand{\baris}[1]{#1 & \tinggi{0.8} & & \\ \hline}
\newcommand{\usulan}[1]{%
\begingroup\renewcommand{\arraystretch}{1.1}
\begin{longtable}{|p{0.96\linewidth}|}
\hline
\textbf{USULAN JUDUL #1} \\ \hline
\bagian{Judul}{(huruf kapital, Arial 14 tebal, paling banyak 20 kata)}{1.4}
\bagian{1.~~Objek/tempat penelitian}{(beserta alamat; bila tanpa objek: Laboratorium Pemrograman/Robotika/Jaringan UM Buton)}{0.9}
\bagian{2.~~Latar belakang masalah}{(jelas, ringkas, sistematis)}{3.0}
\bagian{3.~~Rumusan masalah}{(bernomor)}{1.6}
\bagian{4.~~Tujuan penelitian}{(bernomor, menjawab rumusan masalah)}{1.6}
\end{longtable}
\textbf{Tabel 2.1.} Penelitian terkait\par\vspace{2pt}
\begin{longtable}{|p{0.06\linewidth}|p{0.34\linewidth}|p{0.23\linewidth}|p{0.23\linewidth}|}
\hline
\multicolumn{4}{|p{0.86\linewidth}|}{\textbf{5.~~Penelitian terkait dan daftar pustaka}\ {\footnotesize\itshape (minimal 10 artikel jurnal/prosiding; rujukan gaya APA 7: Penulis, A. B. (Tahun). Judul. Nama Jurnal, 12(1), 45--58. DOI/URL. Kolom Rujukan sekaligus daftar pustaka)}} \\ \hline
\textbf{No.} & \textbf{Rujukan} & \textbf{Persamaan} & \textbf{Perbedaan} \\ \hline
\endhead
\baris{1}\baris{2}\baris{3}\baris{4}\baris{5}\baris{6}\baris{7}\baris{8}\baris{9}\baris{10}
\end{longtable}\endgroup
}

% ---------- Tata letak khusus LaTeX ----------
\IfFontExistsTF{Arial}{\setmainfont{Arial}}{\setmainfont{TeX Gyre Heros}}
\setstretch{1.05}
\setlength{\parindent}{0pt}
\setlength{\parskip}{4pt}
\setlength{\LTpre}{0pt}\setlength{\LTpost}{0pt}
\setlist{nosep,leftmargin=0.6cm}
\renewcommand{\arraystretch}{1.15}
\hyphenpenalty=10000 \exhyphenpenalty=10000 \emergencystretch=3em
\fancypagestyle{plain}{\fancyhf{}\fancyfoot[C]{\thepage}\renewcommand{\headrulewidth}{0pt}}
\pagestyle{plain}

\begin{document}

% ============================ KOP ============================
\begin{tanpagaris}
\begin{longtable}{@{}m{0.17\linewidth}>{\centering\arraybackslash}m{0.60\linewidth}m{0.17\linewidth}@{}}
\logo{2.2cm} &
\textbf{UNIVERSITAS MUHAMMADIYAH BUTON}\par
\textbf{FAKULTAS TEKNIK}\par
\textbf{PROGRAM STUDI REKAYASA SISTEM KOMPUTER} & \\
\end{longtable}
\hrule
\vspace{6pt}
\begin{center}
\textbf{\large\underline{FORMULIR PENGAJUAN JUDUL PENELITIAN}}
\end{center}
\end{tanpagaris}

\begin{tanpagaris}
\begin{longtable}{@{}p{0.34\linewidth}p{0.03\linewidth}p{0.59\linewidth}@{}}
NAMA & : & \titik \\[4pt]
NPM / ANGKATAN & : & \titik \\[4pt]
NO. HP / SUREL & : & \titik \\[4pt]
TANGGAL PENGAJUAN & : & \titik \\
\end{longtable}
\end{tanpagaris}

Dengan ini saya mengajukan tiga usulan judul penelitian/skripsi/tugas akhir. Uraian setiap usulan ada di lembar berikutnya.

\begin{longtable}{|p{0.07\linewidth}|p{0.87\linewidth}|}
\hline
\textbf{No.} & \textbf{Judul (huruf kapital, paling banyak 20 kata)} \\ \hline
1 & \tinggi{1.8} \\ \hline
2 & \tinggi{1.8} \\ \hline
3 & \tinggi{1.8} \\ \hline
\end{longtable}

\begin{kanan}
Baubau, ……………………… 20…

Pengaju,

~

~

~

(\titik)
\end{kanan}

% ============================ USULAN ============================
\begin{halamanbaru}\end{halamanbaru}
\usulan{1}
\begin{halamanbaru}\end{halamanbaru}
\usulan{2}
\begin{halamanbaru}\end{halamanbaru}
\usulan{3}

% ============================ VERIFIKASI ============================
\begin{halamanbaru}\end{halamanbaru}
\begingroup\renewcommand{\arraystretch}{1.3}
\begin{longtable}{|p{0.96\linewidth}|}
\hline
\textbf{HASIL VERIFIKASI} (diisi Ketua Prodi / Dosen Verifikator) \\ \hline
\textbf{Usulan terpilih:}\ \ \cek~Judul 1\ \ \ \cek~Judul 2\ \ \ \cek~Judul 3\ \ \ \cek~Tidak ada \\ \hline
\textbf{Keputusan:}\ \ \cek~Disetujui\ \ \ \cek~Disetujui dengan revisi\ \ \ \cek~Ditolak \\ \hline
\textbf{Catatan/revisi:}\par\tinggi{3.0} \\ \hline
\textbf{Dosen pembimbing} (ditetapkan Ketua Prodi):\newline
Pembimbing I\ \ \ \ : ………………………………………\newline
Pembimbing II\ : ……………………………………… \\ \hline
\end{longtable}\endgroup

\begin{kanan}
Baubau, ……………………… 20…

\textbf{DISETUJUI}\\
\textbf{KETUA PRODI}

~

~

~

\KetuaProdi\\
NIDN. \KetuaProdinidn
\end{kanan}

\end{document}
